\notechapter{N-20220720}{Counting Principles, Binomial Identities, and Mixed Problem Solving}

\section{Counting principles}

The standard counting principles are:
\begin{description}
  \item[Addition principle.] If a task can be done in disjoint cases with \(m_1,m_2,\ldots,m_n\) possibilities, then the total number is \(m_1+\cdots+m_n\).
  \item[Multiplication principle.] If a task is done in stages with \(m_1,m_2,\ldots,m_n\) choices at each stage, then the total number is \(m_1m_2\cdots m_n\).
  \item[Inclusion--exclusion.] For sets \(A_1,\ldots,A_n\), one adds single sizes, subtracts pairwise intersections, adds triple intersections, and so on.

\end{description}

A Venn diagram in the note reviews the three-set formula:
\[
\begin{aligned}
|A\cup B\cup C|&=|A|+|B|+|C|\\
&\quad-|A\cap B|-|A\cap C|-|B\cap C|+|A\cap B\cap C|.
\end{aligned}
\]

\section{A divisibility example}

One problem asks how many positive integers not exceeding \(2021\) are neither multiples of \(6\) nor multiples of \(8\).  The excluded set is
\[
  A\cup B,
\]
where \(A\) is the set of multiples of \(6\), and \(B\) is the set of multiples of \(8\).  Since
\[
  \operatorname{lcm}(6,8)=24,
\]
we have
\[
  |A|=336,\qquad |B|=252,\qquad |A\cap B|=84.
\]
Thus the number excluded is
\[
  336+252-84=504,
\]
and the answer is
\[
  2021-504=1517.
\]
This is exactly what inclusion--exclusion is for: avoid double-counting the common multiples.

\section{Permutations and combinations}

The page distinguishes arrangements and selections:
\[
  A_n^m=\frac{n!}{(n-m)!},\qquad
  C_n^m=\frac{n!}{m!(n-m)!}.
\]
A sequence problem asks for the number of \(10\)-term sequences with \(x+y=10\) and \(x-y=4\), giving \(x=7,y=3\).  Therefore the number of sequences is
\[
  \binom{10}{3}=120.
\]
The thought process is: once the number of one kind of symbol is fixed, the sequence is determined by choosing its positions.

\section{Binomial theorem}

The binomial theorem is
\[
  (a+b)^n=\sum_{k=0}^n\binom nk a^{n-k}b^k.
\]
Several standard identities follow immediately:
\[
  \sum_{k=0}^n\binom nk=2^n,
\]
\[
  \sum_{k=0}^n(-1)^k\binom nk=0,
\]
and by differentiating
\[
  (1+x)^n
\]
and setting \(x=1\),
\[
  \sum_{k=0}^n k\binom nk=n2^{n-1}.
\]
The page writes these identities as examples of how a coefficient formula becomes a counting tool.

\section{Polynomial coefficients}

A worked example asks for the constant term or a particular coefficient in an expansion such as
\[
  (2x-2y)^5.
\]
The term involving \(x^{5-k}y^k\) is
\[
  \binom5k(2x)^{5-k}(-2y)^k.
\]
Thus the coefficient of \(x^2y^3\) comes from \(k=3\):
\[
  \binom53 2^2(-2)^3=-320.
\]
The note's point is that coefficient problems become simple once the target exponent determines \(k\).

\section{Further contest examples}

Counting problems involving subject choices, sign-chart inequalities, and similar-triangle configurations share one skill: case management.  One must decide whether the structure is a sum of cases, a product of stages, or a coefficient extraction.


\section{Binomial identities by stories}

The identity
\[
  \binom n k=\binom{n-1}k+\binom{n-1}{k-1}
\]
comes from choosing a committee of size \(k\) from \(n\) people.  Either a distinguished person is not chosen, or that person is chosen.  This kind of proof is often clearer than algebraic manipulation because it explains what the two terms mean.

\section{When to use multiplication and when to use addition}

The addition principle applies to disjoint alternatives.  The multiplication principle applies to successive independent choices.  Many counting mistakes come from using multiplication when cases overlap, or using addition when a task has stages.  A good habit is to write a sentence: ``first choose ..., then choose ...'' for multiplication, and ``either ..., or ...'' for addition.

\section{A broader reading}

Elementary counting is a first meeting with combinatorial species and generating functions.  Addition corresponds to disjoint union, multiplication corresponds to Cartesian product, and binomial coefficients are the coefficients of a generating function.  This is why the binomial theorem is not just algebra: it is also a counting machine.

\section{Counting principles as operations on combinatorial classes}

The addition and multiplication principles are not merely classroom rules.  They are the first operations on combinatorial classes.  If \(\mathcal A\) and \(\mathcal B\) are disjoint classes, then
\[
  |\mathcal A\sqcup\mathcal B|=|\mathcal A|+|\mathcal B|.
\]
If an object is built by choosing one object from \(\mathcal A\) and one from \(\mathcal B\), then
\[
  |\mathcal A\times\mathcal B|=|\mathcal A||\mathcal B|.
\]
Thus addition corresponds to disjoint union, and multiplication corresponds to Cartesian product.

This viewpoint is useful because it explains why formulas should follow from constructions.  Before computing a number, identify whether the object is a disjoint alternative, a staged product, a quotient by symmetry, or a coefficient in a generating function.

\section{Inclusion--exclusion as Möbius inversion}

The formula
\[
  |A\cup B|=|A|+|B|-|A\cap B|
\]
is a small instance of Möbius inversion on the Boolean lattice.  For a finite family \(A_1,\ldots,A_n\), the number of elements lying in at least one \(A_i\) is
\[
  \left|\bigcup_{i=1}^n A_i\right|
  =\sum_{\varnothing\ne S\subseteq[n]}(-1)^{|S|+1}
  \left|\bigcap_{i\in S}A_i\right|.
\]
The alternating signs are the Möbius function of the subset poset:
\[
  \mu(S,T)=(-1)^{|T|-|S|}.
\]
So inclusion--exclusion is not a trick for Venn diagrams; it is inversion of the operation that replaces exact membership data by cumulative membership data.

\section{Indicator functions and probability language}

For a finite universe \(U\), a set \(A\subseteq U\) has indicator function
\[
  \mathbf 1_A(x)=
  \begin{cases}
  1,&x\in A,\\
  0,&x\notin A.
  \end{cases}
\]
Then
\[
  |A|=\sum_{x\in U}\mathbf 1_A(x).
\]
The two-set inclusion--exclusion identity is the pointwise identity
\[
  \mathbf 1_{A\cup B}=\mathbf 1_A+\mathbf 1_B-\mathbf 1_A\mathbf 1_B.
\]
Summing over \(U\) gives the cardinality formula.

This is also the beginning of probability.  If \(X\) is uniformly chosen from \(U\), then
\[
  \mathbb P(X\in A)=\frac{|A|}{|U|}=\mathbb E[\mathbf 1_A(X)].
\]
Counting and expectation become the same operation after normalization.

\section{Generating functions for binomial identities}

The binomial theorem says
\[
  (1+x)^n=\sum_{k=0}^n\binom nkx^k.
\]
The coefficient of \(x^k\) counts subsets of size \(k\).  Differentiating and setting \(x=1\) gives
\[
  \sum_{k=0}^n k\binom nk=n2^{n-1},
\]
which counts pairs \((S,s)\) where \(S\subseteq[n]\) and \(s\in S\): choose \(s\) first, then choose the rest of \(S\).

This illustrates the right use of generating functions: algebraic operations on series encode operations on counted objects.

\section{Counting up to symmetry: Burnside's lemma}

If a finite group \(G\) acts on a finite set \(X\), the number of orbits is
\[
  |X/G|=\frac1{|G|}\sum_{g\in G}|\operatorname{Fix}(g)|.
\]
This is Burnside's lemma.  It explains why one cannot always divide by the number of symmetries: objects with nontrivial stabilizers have smaller orbits.

Elementary examples include necklaces, colorings of polygons, and arrangements considered up to rotation or reflection.  The advanced lesson is that symmetry changes the unit of counting from elements to orbits.

\section{Counting through structure and symmetry}

The counting principles become powerful when read structurally.  Addition is disjoint union, multiplication is product, inclusion--exclusion is Möbius inversion, binomial coefficients are generating-function coefficients, and symmetry counting is orbit counting.  The elementary skill is to decide which structure the problem is really using before doing arithmetic.
